\chapter{مقدمه و ساختار پروژه}
\label{ch:intro}

\section{زمینه و انگیزه}

امنیت اطلاعات در دهه‌های اخیر به یکی از ارکان اصلی زیرساخت‌های دیجیتال، مخابراتی، مالی، صنعتی و حیاتی تبدیل شده است. بسیاری از سامانه‌های ارتباطی مدرن، از تبادلات بانکی و شبکه‌های داده گرفته تا زیرساخت‌های ابری، سامانه‌های کنترلی و ارتباطات حساس، بر این فرض تکیه دارند که داده‌ها می‌توانند با استفاده از کلیدهای رمزنگاری امن میان دو طرف محافظت و مبادله شوند.
 در رمزنگاری کلاسیک، به‌ویژه در شاخهٔ رمزنگاری کلید عمومی، امنیت بسیاری از سازوکارهای رایج نه از قوانین بنیادین فیزیک، بلکه از دشواری محاسباتی برخی مسائل ریاضی حاصل می‌شود. برای مثال، امنیت الگوریتم‌هایی نظیر \lr{RSA} به دشواری تجزیهٔ اعداد صحیح بزرگ به عوامل اول و امنیت روش‌هایی مانند \lr{ECC} به دشواری حل مسئلهٔ لگاریتم گسسته روی منحنی‌های بیضوی وابسته است.

با وجود موفقیت گستردهٔ این خانواده از الگوریتم‌ها در کاربردهای عملی، این نوع امنیت ذاتاً \textit{امنیت محاسباتی} است، نه \textit{امنیت اطلاعات‌نظری}. به بیان دیگر، این سامانه‌ها تا زمانی امن تلقی می‌شوند که هیچ مهاجم عملی، در بازهٔ زمانی قابل‌قبول، توان محاسباتی کافی برای حل مسائل زیربنایی آن‌ها را در اختیار نداشته باشد. این فرض برای دهه‌ها مبنای طراحی زیرساخت‌های امنیتی بوده است، اما ظهور محاسبات کوانتومی، به‌ویژه الگوریتم شور یا \lr{Shor's algorithm}، نشان داد که یک رایانهٔ کوانتومی مقیاس‌پذیر و دارای تصحیح خطای کافی، از نظر نظری قادر است مسائل فاکتورگیری و لگاریتم گسسته را در زمان چندجمله‌ای حل کند.

تهدید محاسبات کوانتومی تنها به زمان دسترسی عملی به چنین رایانه‌هایی محدود نیست. یکی از نگرانی‌های مهم در امنیت بلندمدت، سناریوی \lr{Harvest Now, Decrypt Later} است؛ به این معنا که مهاجم می‌تواند امروز داده‌های رمزشده را ذخیره کند و در آینده، پس از دستیابی به توان پردازشی لازم، آن‌ها را رمزگشایی کند. بنابراین، حتی اگر رایانه‌های کوانتومی قدرتمند هنوز به‌صورت عملی و گسترده در دسترس نباشند، داده‌هایی که امروز با سازوکارهای سنتی منتقل می‌شوند، ممکن است در آینده محرمانگی بلندمدت خود را از دست بدهند. همین واقعیت، نیاز به روش‌هایی را برجسته می‌کند که امنیت آن‌ها تا حد امکان به توان محاسباتی مهاجم وابسته نباشد.

البته توزیع کلید کوانتومی تنها مسیر پاسخ به تهدید محاسبات کوانتومی نیست. رمزنگاری پساکوانتومی یا \lr{Post-Quantum Cryptography (PQC)} نیز با هدف طراحی الگوریتم‌های کلاسیکی مقاوم در برابر حملات کوانتومی توسعه یافته است. تفاوت اصلی این رویکرد با توزیع کلید کوانتومی در آن است که \lr{PQC} همچنان بر دشواری محاسباتی مسائل ریاضی جدید تکیه دارد، در حالی که توزیع کلید کوانتومی مسئلهٔ توزیع کلید را با استفاده از خواص فیزیکی سامانه‌های کوانتومی بررسی می‌کند. در عمل، این دو رویکرد می‌توانند مکمل یکدیگر باشند؛ به‌ویژه از آن جهت که سامانه‌های \lr{QKD} نیز برای کانال کلاسیک خود به احراز هویت نیاز دارند.

در این بستر، توزیع کلید کوانتومی یا \lr{Quantum Key Distribution (QKD)} به‌عنوان یکی از مهم‌ترین دستاوردهای نظری و عملی در تقاطع فیزیک کوانتومی، نظریهٔ اطلاعات و امنیت ارتباطات مطرح می‌شود. ایدهٔ اصلی \lr{QKD} آن است که دو طرف ارتباطی، که به‌طور قراردادی آلیس (\lr{Alice}) و باب (\lr{Bob}) نامیده می‌شوند، با ارسال حالت‌های کوانتومی و انجام پردازش کلاسیک احراز هویت‌شده، یک کلید محرمانهٔ مشترک تولید کنند. امنیت \lr{QKD} در چارچوب یک مدل امنیتی مشخص و تحت فروض پیاده‌سازی معین تحلیل می‌شود. در این چارچوب، تلاش مهاجم برای کسب اطلاعات از حالت‌های کوانتومی می‌تواند میان اطلاعات مهاجم و اختلال قابل‌مشاهده در داده‌های آلیس و باب مصالحه ایجاد کند. این اختلال از طریق آزمون آماری داده‌های نمونه‌برداری‌شده ارزیابی می‌شود.

این ویژگی، \lr{QKD} را به‌ویژه از سازوکارهای رایج توزیع کلید مبتنی بر رمزنگاری کلید عمومی متمایز می‌کند؛ زیرا امنیت آن، در صورت برقراری فروض مدل امنیتی و وجود یک کانال کلاسیک احراز هویت‌شده، نه بر فرضیات پیچیدگی محاسباتی، بلکه بر اصولی مانند اثر اندازه‌گیری بر حالت کوانتومی، اصل عدم قطعیت و قضیهٔ عدم کلونینگ استوار است. با این حال، این بیان نباید به معنای امنیت مطلق و مستقل از پیاده‌سازی تلقی شود. در سامانه‌های واقعی، اعتبار مدل منبع، آشکارساز، کانال، تصادفی‌سازی، نبود کانال‌های جانبی و نحوهٔ تخصیص پارامترهای امنیتی همگی در نتیجهٔ نهایی اثرگذارند.

در تحلیل‌های مدرن، امنیت کلید نهایی معمولاً با پارامتر شکست \(\epsilon\) بیان می‌شود. به‌طور شهودی، کلید \(\epsilon\)-امن کلیدی است که میزان تمایزپذیری آن از یک کلید ایده‌آلِ یکنواخت و مستقل از اطلاعات مهاجم، با پارامتر \(\epsilon\) کران‌بندی شود. این پارامتر معمولاً احتمال شکست یا عدم تحقق کامل شرایط امنیتی مورد نظر را به‌صورت ترکیبی نمایش می‌دهد. این نگاه به‌ویژه در تحلیل‌های کلید محدود اهمیت دارد، زیرا در سامانه‌های عملی تعداد نمونه‌ها محدود است و نوسانات آماری نمی‌توانند نادیده گرفته شوند.

\section{چالش گذار از \lr{QKD} نظری به \lr{QKD} عملی}

پروتکل \lr{BB84} که در سال ۱۹۸۴ توسط \lr{Bennett} و \lr{Brassard} معرفی شد، نخستین و شناخته‌شده‌ترین پروتکل \lr{QKD} است و همچنان یکی از چارچوب‌های مرجع در تحلیل‌های نظری و پیاده‌سازی‌های عملی به‌شمار می‌رود. در این پروتکل، آلیس بیت‌های اطلاعاتی را در یکی از دو پایهٔ ناسازگار، که در صورت‌بندی ایده‌آل می‌توان آن‌ها را دو پایهٔ متقابلاً نااریب دانست، کدگذاری کرده و برای باب ارسال می‌کند. باب نیز بدون اطلاع از انتخاب پایهٔ آلیس، به‌صورت تصادفی یکی از پایه‌ها را برای اندازه‌گیری انتخاب می‌کند. پس از تبادل کوانتومی، دو طرف از طریق یک کانال کلاسیک احراز هویت‌شده، پایه‌های مورد استفاده را مقایسه کرده و تنها نمونه‌هایی را که در آن‌ها پایه‌ها منطبق بوده‌اند حفظ می‌کنند. این فرایند معمولاً \textit{غربال‌سازی پایه‌ها} یا \lr{basis sifting} نامیده می‌شود. سپس با برآورد نرخ خطا و اجرای مراحل پساپردازش کلاسیک، کلید نهایی استخراج می‌شود.

اهمیت \lr{BB84} از آنجا ناشی می‌شود که اندازه‌گیری حالت کوانتومی در پایه‌ای ناسازگار می‌تواند اختلالی آماری در داده‌ها ایجاد کند. بنابراین، حضور شنودگر نه صرفاً از طریق حدس یا شواهد غیرمستقیم، بلکه از طریق تحلیل آماری نرخ خطا و سایر کمیت‌های مشاهده‌شده ارزیابی می‌شود. در تحلیل امنیت، نرخ خطای مشاهده‌شده تنها یک شاخص عملکردی نیست، بلکه برای کران‌بندی اطلاعات بالقوهٔ مهاجم و، در بسیاری از صورت‌بندی‌ها، برای تخمین کمیت‌هایی مانند خطای فازی نیز به‌کار می‌رود.

با وجود این مبنای نظری، فاصلهٔ میان مدل ایده‌آل و پیاده‌سازی واقعی در عمل تعیین‌کننده است. بسیاری از تحلیل‌های ساده‌شده بر فرض‌هایی مانند وجود منبع تک‌فوتونی کامل، آشکارسازهای بی‌نقص، نبود نویز محیطی و تعداد نامحدود نمونه‌ها تکیه دارند؛ در حالی که هیچ‌یک از این فروض در سامانه‌های واقعی به‌طور کامل برقرار نیست. از همین رو، امنیت و عملکرد یک سامانهٔ \lr{QKD} عملی تنها با بررسی پروتکل ایده‌آل قابل ارزیابی نیست و باید اثر اجزای فیزیکی، آماری و نرم‌افزاری نیز در تحلیل وارد شود.

یکی از محدودیت‌های مهم در پیاده‌سازی عملی \lr{BB84} آن است که ساخت منابع تک‌فوتونی ایده‌آل از نظر فنی دشوار و پرهزینه است. در بسیاری از پیاده‌سازی‌ها، از پالس‌های همدوس تضعیف‌شده یا \lr{Weak Coherent Pulses (WCP)} همراه با تصادفی‌سازی فاز استفاده می‌شود. در این حالت، پالس‌ها را می‌توان به‌صورت مخلوطی آماری از مؤلفه‌های \(n\)-فوتونی با توزیع پواسونی مدل کرد. بنابراین، احتمال غیرصفری برای ارسال پالس‌های چندفوتونی وجود دارد.

وجود پالس‌های چندفوتونی زمینهٔ حملاتی را فراهم می‌کند که در مدل منبع تک‌فوتونی مطرح نیستند. مشهورترین نمونهٔ این تهدید، حملهٔ \lr{Photon-Number Splitting (PNS)} است. در مدل ایده‌آل این حمله، مهاجم می‌تواند با تشخیص تعداد فوتون‌ها، پالس‌های چندفوتونی را شناسایی کند، بخشی از فوتون‌ها را نگه دارد و باقی پالس را به باب برساند. پس از اعلام پایه‌ها، مهاجم می‌تواند از فوتون ذخیره‌شده برای کسب اطلاعات دربارهٔ بخشی از کلید استفاده کند. اگر این رفتار در تحلیل امنیتی لحاظ نشود، ممکن است مهاجم اطلاعاتی دربارهٔ کلید به‌دست آورد، بی‌آنکه نرخ خطای مشاهده‌شده به‌تنهایی برای آشکارسازی حمله کافی باشد. بنابراین، کنترل اثر پالس‌های چندفوتونی بخش مهمی از تحلیل امنیت عملی سامانه‌های مبتنی بر \lr{WCP} است.

برای مقابله با این ضعف ساختاری، روش حالت‌های فریب یا \lr{Decoy-State Method} معرفی شد که از مهم‌ترین پیشرفت‌ها در عملی‌سازی \lr{QKD} محسوب می‌شود. در این رویکرد، آلیس به‌جای ارسال تمام پالس‌ها با یک شدت ثابت، به‌صورت تصادفی از چند شدت متفاوت استفاده می‌کند؛ معمولاً یک شدت سیگنال، یک یا چند شدت فریب ضعیف، و در بسیاری از موارد یک حالت خلأ یا \lr{vacuum decoy}. هدف این روش آن است که رفتار کانال و آشکارسازی برای مؤلفه‌های فوتونی مختلف به‌صورت آماری تخمین زده شود.

در روش حالت‌های فریب، با فرض تصادفی‌سازی کامل فاز و نبود کانال جانبی وابسته به شدت، حالت‌های \(n\)-فوتونی حاصل از شدت‌های مختلف از دید مهاجم غیرقابل‌تمایز در نظر گرفته می‌شوند. در نتیجه، بازده و نرخ خطای متناظر با هر مؤلفهٔ \(n\)-فوتونی باید میان حالت‌های سیگنال و فریب مشترک باشد. آلیس و باب با مقایسهٔ آمار شدت‌های مختلف، می‌توانند بر کمیت‌هایی مانند بازده تک‌فوتونی \(Y_1\) و نرخ خطای تک‌فوتونی \(e_1\) کران بگذارند. این تخمین‌ها زیربنای استخراج نرخ کلید امن در حضور منابع غیرایده‌آل را تشکیل می‌دهند.

به بیان دیگر، روش حالت‌های فریب امکان می‌دهد سهم پالس‌های تک‌فوتونی، که برای امنیت پروتکل نقش اصلی دارند، از سهم پالس‌های چندفوتونی، که بالقوه آسیب‌پذیرند، به‌صورت آماری تفکیک شود. در غیاب این روش، عملکرد بسیاری از سامانه‌های مبتنی بر \lr{WCP} در فواصل متوسط و بلند یا از نظر امنیتی ناکافی می‌بود یا نرخ کلید آن‌ها به‌شدت افت می‌کرد. با بهره‌گیری از این تکنیک، می‌توان در شرایط واقع‌گرایانه و با در نظر گرفتن محدودیت‌های عملی، به نرخ کلید امن قابل‌استفاده دست یافت.

افزون بر مدل منبع، تحلیل سامانه‌های \lr{QKD} باید میان دو رژیم مهم تمایز قائل شود: رژیم \textit{مجانبی} یا \lr{asymptotic} و رژیم \textit{کلید محدود} یا \lr{finite-key}. در مدل مجانبی فرض می‌شود که تعداد پالس‌های مبادله‌شده آن‌قدر زیاد است که نوسانات آماری ناچیز باشند و پارامترها تقریباً دقیق تخمین زده شوند. اما در عمل، هر آزمایش و هر سامانهٔ عملی با تعداد محدودی از پالس‌ها و زمان اجرایی محدود روبه‌رو است. در این وضعیت، نوسانات آماری و خطاهای نمونه‌برداری می‌توانند اثر مهمی بر برآورد کمیت‌های امنیتی بگذارند.

از این رو، تحلیل‌های کلید محدود برای ارزیابی امنیت عملی سامانه‌هایی با اندازهٔ بلوک محدود ضروری‌اند. تحلیل مجانبی همچنان برای فهم رفتار پایه، مقایسه‌های اولیه و اعتبارسنجی تحلیلی ارزشمند است، اما اثر نوسانات آماری و بودجهٔ امنیتی را به‌طور کامل بازنمایی نمی‌کند. در بسیاری از صورت‌بندی‌های کلید محدود، طول کلید نهایی به کمیت‌هایی مانند کران پایین تعداد رویدادهای تک‌فوتونی، کران بالای خطای متناظر، نشت اطلاعات در تصحیح خطا و جمله‌های وابسته به پارامترهای امنیتی بستگی دارد. به‌صورت مفهومی، بیشینهٔ طول کلید قابل استخراج را می‌توان با رابطه‌ای از شکل زیر کران‌بندی کرد:
\[
\ell \lesssim s_{1}^{L}\bigl[1-h_2(e_{1}^{U})\bigr]
-\lambda_{\mathrm{EC}}
-\Delta(\epsilon),
\]
در این رابطه، \(h_2(\cdot)\) تابع آنتروپی دودویی، \(s_{1}^{L}\) کران پایین سهم تک‌فوتونی‌ها، \(e_{1}^{U}\) کران بالای نرخ خطای متناظر، \(\lambda_{\mathrm{EC}}\) نشت اطلاعات در مرحلهٔ تصحیح خطا، و \(\Delta(\epsilon)\) جمله‌ای وابسته به بودجهٔ امنیتی است. شکل دقیق این رابطه به اثبات امنیتی و فروض مدل وابسته است.

بنابراین، مطالعه و توسعهٔ سامانه‌های \lr{QKD} تنها به شناخت اصول بنیادی فیزیک کوانتومی محدود نمی‌شود. چنین مطالعه‌ای مستلزم درک تعامل میان منبع نوری، کانال انتقال، نویزها، آشکارسازها، فرایندهای آماری، اثبات‌های امنیتی و مراحل پساپردازش کلاسیک است. پروژه‌ای که در این گزارش معرفی می‌شود، در همین نقطه قرار دارد: طراحی و پیاده‌سازی یک بستر شبیه‌سازی که بتواند فاصلهٔ میان مدل‌های نظری و رفتار عملی سامانه‌های \lr{QKD} را به‌صورت کنترل‌شده و قابل‌اعتبارسنجی مطالعه کند و امکان تحلیل، اعتبارسنجی، مقایسه و توسعهٔ سناریوهای پیشرفته را فراهم کند.

\section{ضرورت شبیه‌سازی در سامانه‌های \lr{QKD}}


هرچند اصول بنیادی \lr{QKD} در چارچوب نظری به‌خوبی شناخته شده‌اند، گذار از اثبات امنیت روی کاغذ به یک سامانهٔ عملی و قابل‌اتکا، فرایندی پیچیده و چندلایه است. در یک پیاده‌سازی واقعی، هیچ‌یک از اجزای فیزیکی رفتار کاملاً ایده‌آل ندارند و همین انحراف از ایده‌آل، هم بر عملکرد و هم بر امنیت نهایی اثر می‌گذارد. در نتیجه، طراحی، تحلیل و بهینه‌سازی سامانه‌های \lr{QKD} بدون اتکا به ابزارهای شبیه‌سازی دقیق، در بسیاری از موارد ناقص خواهد بود.

در لایهٔ منبع، پالس‌های نوری می‌توانند با نوساناتی در شدت، فاز، زمان‌بندی و حتی توزیع آماری تعداد فوتون همراه باشند. در عمل، منبع ممکن است دچار جیتر شدت یا \lr{intensity jitter} شود، پالس‌ها از نظر شکل زمانی یکنواخت نباشند، یا در برخی شرایط از مدل ایده‌آل فاصله بگیرند. این انحرافات می‌توانند مستقیماً بر تخمین پارامترهای امنیتی و همچنین بر نرخ کلید مؤثر اثر بگذارند. از سوی دیگر، کانال انتقال، به‌ویژه در بستر فیبر نوری، با تضعیف، پراکندگی، وابستگی طول‌موجی، ناپایداری‌های محیطی و نویز پس‌زمینه همراه است. در فواصل بلند، کاهش بازده انتقال می‌تواند بر نرخ آشکارسازی و در نتیجه بر امکان‌پذیری تولید کلید امن اثر قابل‌ملاحظه‌ای بگذارد.

در سمت گیرنده نیز وضعیت به همین اندازه پیچیده است. آشکارسازهای تک‌فوتونی، به‌خصوص آشکارسازهای مبتنی بر \lr{SPAD}، دارای رفتارهایی هستند که در تحلیل ایده‌آل معمولاً نادیده گرفته می‌شوند اما در عمل نقش تعیین‌کننده دارند. از جملهٔ این رفتارها می‌توان به شمارش تاریک یا \lr{dark count}، زمان مرده یا \lr{dead time}، پس‌تپش یا \lr{afterpulsing}، محدودیت‌های گیتینگ زمانی، اشباع‌شدگی در نرخ‌های بالا، و سیاست‌های متفاوت در برخورد با رخدادهای هم‌زمان آشکارسازها یا \lr{double clicks} اشاره کرد.

هر یک از این عوامل می‌تواند هم بر نرخ خطای کوانتومی یا \lr{QBER} و هم بر نرخ کلید نهایی اثر بگذارد. حتی در برخی موارد، عدم مدل‌سازی دقیق همین جزئیات موجب می‌شود نتیجهٔ شبیه‌سازی به‌ظاهر امیدوارکننده باشد، در حالی که سامانهٔ واقعی آن عملکرد را محقق نکند. بنابراین، شبیه‌سازی زمانی ارزشمند است که بتواند اثر این جزئیات را با دقت مناسب و در چارچوب فروض روشن وارد تحلیل کند.

اهمیت موضوع زمانی بیشتر روشن می‌شود که بدانیم بسیاری از تصمیمات طراحی در \lr{QKD} به‌شدت پارامترمحور هستند. انتخاب شدت‌های سیگنال و فریب، احتمال استفاده از هر شدت، نرخ تکرار پالس، بازهٔ زمانی گیت، مدل تلفات کانال، بازده آشکارسازی، و سیاست پردازش داده در لایهٔ کلاسیک، همگی بر خروجی نهایی اثر می‌گذارند. از آنجا که این پارامترها معمولاً به‌صورت کوپل‌شده عمل می‌کنند، فهم اثر مستقل هر یک بدون وجود یک محیط شبیه‌سازی کنترل‌شده دشوار است. برای مثال، افزایش شدت سیگنال ممکن است نرخ آشکارسازی را بالا ببرد، اما هم‌زمان احتمال رخداد پالس‌های چندفوتونی را نیز افزایش دهد. از سوی دیگر، کاهش بیش از حد شدت، اگرچه آسیب‌پذیری در برابر حملات مرتبط با چندفوتونی‌ها را کاهش می‌دهد، ممکن است به افت شدید نرخ آشکارسازی و غلبهٔ نویز تاریک منجر شود. چنین مصالحه‌هایی تنها از طریق تحلیل کمی و تکرارپذیر قابل ارزیابی‌اند.

از منظر پژوهشی و مهندسی، آزمایش‌های سخت‌افزاری \lr{QKD} معمولاً پرهزینه، زمان‌بر، حساس به شرایط محیطی، و محدود از نظر قابلیت پیمایش فضای پارامتر هستند. ساخت یا پیکربندی مجدد یک آزمایش برای هر مجموعهٔ جدید از پارامترها، نه‌تنها هزینهٔ زمانی و مالی دارد، بلکه در بسیاری از موارد امکان اجرای گستردهٔ سناریوهای مقایسه‌ای را محدود می‌کند. در مقابل، یک شبیه‌ساز می‌تواند تعداد زیادی سناریو را در یک محیط کنترل‌شده بررسی کند، نقاط بحرانی عملکرد را بیابد، فرضیه‌ها را پیش از پیاده‌سازی سخت‌افزاری آزمایش کند، و مسیر طراحی آزمایش واقعی را هدایت نماید. بنابراین، در بسیاری از مطالعات عملی \lr{QKD}، شبیه‌سازی نه‌تنها ابزاری کمکی، بلکه بخشی از فرایند طراحی، اعتبارسنجی و بهینه‌سازی سامانه محسوب می‌شود.

یک شبیه‌ساز مفید برای \lr{QKD} نباید تنها به محاسبهٔ روابط بسته و فرمول‌های سادهٔ مجانبی محدود شود. بسیاری از تحلیل‌های نظری در ساده‌ترین حالت، پارامترهایی مانند بهرهٔ آشکارسازی، نرخ خطا یا نرخ کلید را بر اساس روابط تحلیلی در شرایط ایده‌آل یا نیمه‌واقعی برآورد می‌کنند. این روابط برای فهم شهودی و اعتبارسنجی اولیه ارزشمندند، اما برای مدل‌سازی رفتار واقعی سامانه، به‌خصوص در مقیاس کلید محدود، کافی نیستند. در عمل، رویدادها ماهیتاً گسسته و آماری‌اند: فوتون یا به آشکارساز می‌رسد یا نمی‌رسد، کلیک رخ می‌دهد یا رخ نمی‌دهد، پالس ممکن است به‌دلیل زمان مرده از دست برود، و نوسانات تصادفی می‌توانند بر تخمین پارامترها اثر بگذارند.

بنابراین، شبیه‌سازی رویدادمحور و مبتنی بر مونت‌کارلو ابزاری مهم برای بازنمایی رخدادهای گسسته، وابستگی‌های زمانی و نوسانات آماری است. با این حال، نتایج چنین شبیه‌سازی‌ای باید در سناریوهای مناسب با روابط تحلیلی یا بنچمارک‌های مستقل مقایسه شوند تا هزینهٔ محاسباتی بیشتر آن به افزایش قابل‌سنجش در دقت مدل منجر شود.

علاوه بر این، در یک سامانهٔ پژوهشی، صرف شبیه‌سازی لایهٔ کوانتومی کافی نیست. پس از مرحلهٔ آشکارسازی، مجموعه‌ای از پردازش‌های کلاسیک شامل غربال‌سازی پایه‌ها، نمونه‌برداری برای برآورد خطا، تخمین پارامترهای حالت فریب، مدل‌سازی هزینهٔ تصحیح خطا، تخصیص بودجه‌های امنیتی، و محاسبهٔ طول یا نرخ کلید امن انجام می‌شود. اگر این مراحل با دقت آماری کافی مدل نشوند، نتیجهٔ شبیه‌سازی از منظر امنیتی قابل‌استناد نخواهد بود. به بیان دیگر، شبیه‌ساز مطلوب باید میان لایهٔ فیزیکی سامانه و لایهٔ نظریهٔ اطلاعات ارتباط برقرار کند.

در نهایت، شبیه‌سازی دقیق امکان می‌دهد مدل‌های فیزیکی مختلف، اثبات‌های امنیتی متفاوت، یا معماری‌های اجرایی متنوع در یک بستر واحد با یکدیگر مقایسه شوند. چنین قابلیتی برای پژوهش‌هایی که هدف آن‌ها توسعهٔ مدل جدید، اعتبارسنجی یک تقریب فیزیکی، سنجش اثر یک پارامتر، یا مقایسهٔ چند روش تحلیل امنیت است، اهمیت زیادی دارد. از همین رو، توسعهٔ یک بستر شبیه‌سازی توسعه‌پذیر و اعتبارسنجی‌پذیر، نه‌تنها برای پاسخ‌گویی به نیازهای این پروژه، بلکه به‌عنوان زیرساختی برای مطالعات بعدی نیز ارزشمند است.

\section{بیان مسئله و هدف پروژه}

مسئلهٔ اصلی این پروژه، طراحی و پیاده‌سازی یک فریم‌ورک شبیه‌سازی رویدادمحور برای تحلیل سامانه‌های \lr{DV-QKD} است که بتواند مدل‌سازی فیزیکی، رفتار آماری، محدودیت‌های سخت‌افزاری، تحلیل‌های امنیتی و بخشی از نیازهای معماری شبکه‌ای را در یک بستر نرم‌افزاری واحد ترکیب کند. تمرکز نسخهٔ فعلی فریم‌ورک بر خانوادهٔ سامانه‌های مبتنی بر \lr{BB84}، منابع مبتنی بر پالس‌های همدوس تضعیف‌شده، روش حالت‌های فریب و تحلیل‌های مجانبی و کلید محدود است. با این حال، معماری آن به‌گونه‌ای طراحی شده است که افزودن مدل‌ها و سناریوهای جدید در آینده با کمترین وابستگی غیرضروری به بخش‌های موجود انجام شود.

در این گزارش، منظور از فریم‌ورک، مجموعه‌ای از ماژول‌های نرم‌افزاری، رابط‌های برنامه‌نویسی، ابزارهای اجرا، و اسکریپت‌های اعتبارسنجی است که به‌صورت یکپارچه برای تعریف، اجرا و تحلیل سناریوهای \lr{QKD} به‌کار می‌روند. هدف این فریم‌ورک تنها محاسبهٔ چند خروجی نهایی نیست، بلکه بازنمایی مرحله‌به‌مرحلهٔ زنجیرهٔ تولید پالس، عبور از کانال، اثر نویزها، آشکارسازی، تجمیع آمار، تخمین پارامترها و محاسبهٔ کمیت‌های امنیتی است.

هدف‌های اصلی پروژه را می‌توان به‌صورت زیر خلاصه کرد:

\begin{enumerate}
	\item فراهم‌کردن محیطی برای مدل‌سازی پدیده‌های فیزیکی مؤثر بر عملکرد سامانه‌های \lr{QKD}، از جمله منبع، کانال، نویز و آشکارساز؛
	
	\item ایجاد امکان مقایسهٔ چند خانواده از تحلیل‌های امنیتی تحت شرایط فیزیکی و آماری مشترک؛
	
	\item پشتیبانی از تحلیل‌های کلید محدود، در کنار تحلیل‌های مجانبی، برای بررسی اثر نوسانات آماری و بودجه‌های امنیتی؛
	
	\item فراهم‌کردن زیرساخت لازم برای اجرای شبیه‌سازی‌های رویدادمحور در مقیاس بزرگ و با قابلیت پردازش دسته‌ای یا موازی؛
	
	\item فراهم‌سازی معماری نرم‌افزاری ماژولار برای افزودن مدل‌های جدید منبع، کانال، آشکارساز، حمله یا اثبات امنیت؛
	
	\item ایجاد قابلیت اولیه برای مطالعهٔ سناریوهای چندگرهی و شبکه‌ای مبتنی بر لینک‌های \lr{QKD}.
\end{enumerate}

در نسخهٔ فعلی، تصحیح خطا و فشرده‌سازی محرمانگی در سطح اثر آن‌ها بر طول یا نرخ کلید نهایی مدل می‌شوند. به بیان دیگر، تمرکز اصلی فریم‌ورک بر محاسبهٔ کمیت‌های امنیتی و هزینه‌های اطلاعاتی متناظر با مراحل پساپردازش است، نه الزاماً اجرای کامل یک پروتکل عملی تصحیح خطا یا اعمال فیزیکی یک تابع هش جهانی برای تولید بیت‌های نهایی کلید. جزئیات دقیق این مدل‌سازی در فصل‌های مربوط به پساپردازش و تحلیل امنیت ارائه می‌شود.

همچنین، اگرچه تمرکز اصلی فریم‌ورک بر تحلیل لینک‌های \lr{DV-QKD} است، معماری آن به‌گونه‌ای توسعه یافته که بتوان چندین لینک را در قالب سناریوهای شبکه‌ای مبتنی بر رلهٔ مورد اعتماد ترکیب کرد. در چنین معماری‌ای، امنیت انتهابه‌انتها به اعتماد به گره‌های میانی وابسته است؛ زیرا این گره‌ها در فرایند انتقال یا بازتوزیع کلید به اطلاعات کلیدی دسترسی دارند.

\section{نمای کلی فریم‌ورک شبیه‌سازی}

پروژهٔ حاضر یک فریم‌ورک رویدادمحور برای شبیه‌سازی سامانه‌های \lr{Discrete-Variable Quantum Key Distribution (DV-QKD)} ارائه می‌کند. این فریم‌ورک با هدف ترکیب مدل‌سازی فیزیکی، تحلیل آماری، محاسبهٔ کمیت‌های امنیتی و معماری نرم‌افزاری توسعه‌پذیر طراحی شده است. رویکرد محاسباتی آن بر پایهٔ شبیه‌سازی مونت‌کارلو و پردازش رویدادهاست؛ به این معنا که پالس‌ها، آشکارسازی‌ها و رخدادهای مرتبط با آن‌ها به‌صورت نمونه‌های گسسته تولید و تحلیل می‌شوند.

در این معماری، هر پالس نوری به‌صورت یک رویداد مستقل با شناسه و ویژگی‌های مشخص مدل می‌شود. ویژگی‌های پالس می‌توانند در لایه‌های مختلف سامانه، از لحظهٔ تولید در منبع نوری تا عبور از کانال، مواجهه با تلفات و نویز، رسیدن یا نرسیدن به آشکارساز، ثبت کلیک، و ورود به مراحل پساپردازش کلاسیک، تغییر کنند. این رویکرد امکان می‌دهد نه‌تنها خروجی نهایی، بلکه مسیر ایجاد آن و سهم هر بخش از سامانه نیز تحلیل شود.

زنجیرهٔ شبیه‌سازی را می‌توان به‌صورت مفهومی در چند لایه توصیف کرد:

\begin{enumerate}
	\item \textbf{لایهٔ منبع:} منبع نوری پالس‌ها را بر اساس مدل آماری انتخاب‌شده تولید می‌کند. در سناریوهای مبتنی بر منابع همدوس تضعیف‌شده، تعداد فوتون‌ها معمولاً با توزیع پواسونی و شدت‌های مختلف حالت فریب مدل می‌شود.
	
	\item \textbf{لایهٔ کانال:} کانال انتقال اثر تلفات، طول مسیر، بازده انتقال و در صورت نیاز نویزهای پس‌زمینه را بر پالس‌ها اعمال می‌کند.
	
	\item \textbf{لایهٔ مدل تهدید:} در صورت تعریف سناریوی حمله یا مدل تهدید، اثر آن با فروض مشخص به فرایند شبیه‌سازی افزوده می‌شود.
	
	\item \textbf{لایهٔ آشکارسازی:} زیرسامانهٔ آشکارسازی رفتارهایی مانند احتمال کلیک، شمارش تاریک، زمان مرده، گیتینگ زمانی، رخدادهای هم‌زمان و سیاست رسیدگی به آن‌ها را مدل می‌کند.
	
	\item \textbf{لایهٔ پساپردازش کلاسیک:} داده‌های خام وارد مراحل غربال‌سازی، تجمیع شمارش‌ها، تخمین پارامترهای امنیتی، اعمال مدل تصحیح خطا و محاسبهٔ نرخ یا طول کلید امن می‌شوند.
	
	\item \textbf{لایهٔ شبکه:} در صورت استفاده از سناریوهای چندگرهی، خروجی لینک‌ها در سطح شبکه ترکیب می‌شود و مفاهیمی مانند استخر کلید، مسیرهای جایگزین و رمزنگاری \lr{hop-by-hop} قابل بررسی می‌شوند.
\end{enumerate}

طراحی نرم‌افزاری فریم‌ورک بر چند اصل استوار است. نخست، منطق مربوط به پارامترها، منبع، کانال، آشکارساز، تحلیل امنیت، بهینه‌سازی و شبکه تا حد امکان به‌صورت ماژول‌های مستقل سازمان‌دهی شده است. این تفکیک باعث می‌شود تغییر یا جایگزینی یک مدل، کمترین اثر را بر سایر بخش‌ها داشته باشد. دوم، هویت و ویژگی‌های رویدادها در طول پردازش حفظ می‌شود تا تحلیل خطا، اشکال‌زدایی و اعتبارسنجی ممکن باشد. سوم، رابط‌های نرم‌افزاری به‌گونه‌ای طراحی شده‌اند که افزودن مدل‌های جدید بدون بازنویسی کل سیستم انجام شود. چهارم، خروجی‌ها باید بتوانند در سناریوهای ساده‌تر با روابط تحلیلی شناخته‌شده یا بنچمارک‌های مستقل مقایسه شوند.

با توجه به حجم بالای رویدادها در شبیه‌سازی‌های بزرگ، فریم‌ورک از پردازش دسته‌ای و سازوکارهای توزیع بار محاسباتی نیز پشتیبانی می‌کند. در چنین حالتی، حفظ یکپارچگی شناسهٔ پالس‌ها، همگام‌سازی نتایج دسته‌ای و مدیریت رویدادهایی که اثرشان از یک دستهٔ پردازشی به دستهٔ بعد منتقل می‌شود، اهمیت زیادی دارد. ارزیابی کمی کارایی، شامل زمان اجرا، مصرف حافظه و رفتار مقیاس‌پذیری، در فصل مربوط به نتایج و اعتبارسنجی ارائه می‌شود.

در بخش تحلیل امنیت، فریم‌ورک امکان مقایسهٔ چند خانواده از روش‌های محاسبهٔ نرخ کلید را فراهم می‌کند؛ از جمله تحلیل‌های مجانبی کلاسیک برای \lr{Decoy-State BB84}، تحلیل‌های کلید محدود مبتنی بر کران‌های آماری، و نسخه‌هایی که با هدف ارائهٔ کران‌های عددی تنگ‌تر توسعه یافته‌اند. هدف از قرار دادن این روش‌ها در یک بستر مشترک آن است که اثر تفاوت اثبات امنیتی از اثر تفاوت مدل فیزیکی یا فروض نرم‌افزاری تا حد امکان جدا شود.

بهینه‌سازی پارامترها نیز یکی از قابلیت‌های مهم فریم‌ورک است. انتخاب شدت سیگنال \(\mu\)، شدت فریب \(\nu\)، احتمال استفاده از هر شدت، و برخی تنظیمات کانال یا آشکارساز، اثر چشمگیری بر عملکرد سامانه دارد. از این رو، فریم‌ورک می‌تواند از روش‌های عددی و حلگرهای بهینه‌سازی برای پیمایش فضای پارامتر و شناسایی پیکربندی‌های مناسب‌تر استفاده کند. در مسائل خطی یا مراحل کران‌بندی مرتبط با تخمین پارامترها نیز استفاده از ابزارهای برنامه‌ریزی خطی قابل بررسی است.

در مجموع، این فریم‌ورک را می‌توان یک زیرساخت پژوهشی برای مطالعهٔ سامانه‌های \lr{QKD} دانست که با هدف توسعهٔ تدریجی، مقایسهٔ روش‌ها، تولید نتایج تکرارپذیر و پشتیبانی از مطالعات پیچیده‌تر در حوزهٔ \lr{QKD} و شبکه‌های کوانتومی طراحی شده است.

\section{مشارکت‌ها و قابلیت‌های اصلی فریم‌ورک}

مشارکت‌های اصلی این پروژه را می‌توان در چند محور خلاصه کرد:

\begin{enumerate}
	\item \textbf{شبیه‌سازی رویدادمحور لینک QKD:} توسعهٔ یک زنجیرهٔ شبیه‌سازی که پالس‌ها، کانال، نویز، آشکارسازی و پساپردازش را به‌صورت مرحله‌به‌مرحله مدل می‌کند.
	
	\item \textbf{پشتیبانی از منابع همدوس تضعیف‌شده و حالت‌های فریب:} مدل‌سازی منابع مبتنی بر \lr{WCP} و استفاده از شدت‌های مختلف برای تخمین آماری کمیت‌های تک‌فوتونی.
	
	\item \textbf{در نظر گرفتن محدودیت‌های سخت‌افزاری:} وارد کردن عواملی مانند تلفات کانال، شمارش تاریک، زمان مرده، گیتینگ و رخدادهای هم‌زمان آشکارسازها در تحلیل عملکرد.
	
	\item \textbf{ترکیب تحلیل مجانبی و کلید محدود:} فراهم‌کردن امکان بررسی سامانه در دو رژیم مجانبی و کلید محدود، با توجه به اثر نوسانات آماری و بودجهٔ امنیتی.
	
	\item \textbf{معماری ماژولار و توسعه‌پذیر:} تفکیک لایه‌های منبع، کانال، نویز، آشکارساز، اثبات امنیت، بهینه‌سازی و شبکه برای تسهیل نگهداری و افزودن مدل‌های جدید.
	
	\item \textbf{پشتیبانی از اعتبارسنجی داخلی:} استفاده از آزمون‌ها و سناریوهای کنترلی برای بررسی سازگاری روابط فیزیکی، پایداری خروجی‌ها و اتصال صحیح لایه‌های نرم‌افزاری.
	
	\item \textbf{زیرساخت اولیه برای سناریوهای شبکه‌ای:} فراهم‌کردن امکان مطالعهٔ چند لینک و چند گره در قالب معماری‌های مبتنی بر رلهٔ مورد اعتماد.
\end{enumerate}

این موارد دامنهٔ اصلی پروژه را مشخص می‌کنند و مرز میان هدف پژوهشی، مدل‌سازی فیزیکی و پیاده‌سازی نرم‌افزاری را روشن می‌سازند. جزئیات هر یک از این محورها در فصل‌های بعدی گزارش ارائه خواهد شد.

\section{ساختار نرم‌افزاری و اجرای پروژه}

ساختار نرم‌افزاری پروژه بر پایهٔ تفکیک مسئولیت‌ها طراحی شده است. در سطح مفهومی، ماژول‌های اصلی فریم‌ورک شامل مدیریت پارامترها، تعریف ساختارهای داده، مدل‌سازی منبع، کانال و نویز، مدل‌سازی آشکارساز، تحلیل امنیت، بهینه‌سازی، اجرای سناریوها و آزمون‌های اعتبارسنجی هستند. این تفکیک باعث می‌شود که هر بخش از سامانه بتواند مستقل‌تر توسعه یابد و اثر تغییرات آن بر سایر بخش‌ها کنترل شود.

در ریشهٔ پروژه، فایل‌های اجرایی و اسکریپت‌های هماهنگ‌کننده قرار دارند. این فایل‌ها نقش نقطهٔ ورود برای اجرای سناریوهای اصلی، اجرای پردازش‌های دسته‌ای یا موازی، اعتبارسنجی روابط فیزیکی و اجرای نمونه‌سناریوهای شبکه‌ای را بر عهده دارند. جزئیات نام فایل‌ها، وابستگی‌های دقیق آن‌ها و نحوهٔ اجرای هر سناریو در فصل پیاده‌سازی یا پیوست راهنمای اجرا ارائه می‌شود.

هستهٔ فریم‌ورک در قالب مجموعه‌ای از ماژول‌ها سازمان‌دهی شده است. ماژول مدیریت پارامترها، تنظیمات فیزیکی، امنیتی و اجرایی سناریو را نگهداری می‌کند. ماژول‌های داده‌ای، ساختارهای لازم برای انتقال شمارش‌ها، رویدادها و نتایج میانی میان لایه‌ها را تعریف می‌کنند. ماژول‌های فیزیکی، روابط مربوط به تلفات، تولید فوتون، نویز و آشکارسازی را پیاده‌سازی می‌کنند. لایهٔ اثبات‌های امنیتی نیز از داده‌های خروجی لایهٔ فیزیکی استفاده می‌کند تا کمیت‌هایی مانند نرخ کلید، طول کلید، کران‌های آماری و اثر هزینه‌های پساپردازش را محاسبه کند.

برای کنترل صحت توسعه، مجموعه‌ای از آزمون‌های واحد، آزمون‌های یکپارچه‌سازی و در صورت نیاز آزمون‌های رگرسیون در پروژه در نظر گرفته شده است. آزمون‌های واحد برای سنجش رفتار توابع یا کلاس‌های کوچک و مستقل به‌کار می‌روند؛ برای مثال، بررسی صحت محاسبهٔ تلفات یا رفتار یک ساختار داده در شرایط مرزی. آزمون‌های یکپارچه‌سازی تعامل میان ماژول‌ها را بررسی می‌کنند؛ برای نمونه، سازگاری خروجی آشکارساز با ورودی تحلیل امنیت، یا حفظ هویت رویدادها در پردازش دسته‌ای. آزمون‌های رگرسیون نیز برای جلوگیری از بازگشت خطاهای قدیمی پس از اعمال تغییرات جدید مفیدند.

از نظر اجرایی، پیاده‌سازی اصلی پروژه با زبان پایتون انجام شده است. انتخاب پایتون به دلیل اکوسیستم علمی غنی آن، سرعت بالای توسعه، سهولت اعتبارسنجی پژوهشی و دسترسی به کتابخانه‌های عددی و آماری انجام شده است. کتابخانه‌هایی مانند \texttt{NumPy}، \texttt{SciPy}، \texttt{pandas} و \texttt{pytest} در بخش‌های مختلف پروژه برای محاسبات عددی، بهینه‌سازی، مدیریت داده‌های خروجی و اجرای آزمون‌ها مورد استفاده قرار می‌گیرند. بسته به تنظیمات پروژه و ویژگی‌های فعال‌شده، ممکن است وابستگی‌های کمکی دیگری نیز برای ثبت گزارش، ترسیم نمودار یا مدیریت داده‌ها استفاده شوند.

برای اجرای یک شبیه‌سازی پایه، کاربر معمولاً پارامترهای سناریو را در پیکربندی مربوط تنظیم کرده و سپس نقطهٔ ورود مناسب را اجرا می‌کند. در طی اجرا، سیستم پارامترها را بارگذاری کرده، اجزای لازم را نمونه‌سازی می‌کند، رویدادها را طبق سناریوی مشخص تولید و پردازش می‌نماید، و در نهایت خروجی‌ها را در قالب فایل‌های ساخت‌یافته ذخیره می‌کند. در طراحی خروجی‌ها، پایداری داده‌ها و کاهش خطر خرابی فایل در اجراهای طولانی‌مدت یا دسته‌ای مورد توجه قرار گرفته است.

جزئیات کامل ساختار فایل‌ها، نام اسکریپت‌های اجرایی، وابستگی‌های نرم‌افزاری، نحوهٔ نصب، روش اجرای سناریوهای پایه و پیشرفته، و شیوهٔ بازتولید نتایج در فصل پیاده‌سازی و پیوست راهنمای اجرا ارائه خواهد شد. به این ترتیب، فصل حاضر از ورود به جزئیات پایین‌سطح نرم‌افزاری پرهیز می‌کند و تمرکز خود را بر معرفی مسئله، دامنهٔ پروژه و معماری مفهومی فریم‌ورک نگه می‌دارد.

\section{ساختار گزارش}

ادامهٔ این گزارش به‌گونه‌ای سازمان‌دهی شده است که خواننده ابتدا با مبانی نظری و سپس با مدل‌سازی، پیاده‌سازی، اعتبارسنجی و نتایج پروژه آشنا شود. به‌طور کلی، ساختار فصل‌های بعدی به شرح زیر است:

\begin{itemize}
	\item در فصل مبانی نظری، اصول \lr{QKD}، پروتکل \lr{BB84}، منابع همدوس تضعیف‌شده، روش حالت‌های فریب، تحلیل مجانبی و تحلیل کلید محدود با جزئیات بیشتری بررسی می‌شوند.
	
	\item در فصل مدل فیزیکی و آماری، مدل‌های مورد استفاده برای منبع، کانال، نویز، آشکارساز و رخدادهای رویدادمحور معرفی می‌شوند.
	
	\item در فصل تحلیل امنیت، روش‌های محاسبهٔ نرخ یا طول کلید امن، کران‌بندی پارامترها، اثر نوسانات آماری و نقش بودجهٔ امنیتی توضیح داده می‌شوند.
	
	\item در فصل معماری و پیاده‌سازی نرم‌افزار، ساختار ماژول‌ها، ارتباط میان لایه‌ها، نحوهٔ مدیریت داده‌ها، پردازش دسته‌ای و جزئیات اجرایی فریم‌ورک ارائه می‌شود.
	
	\item در فصل اعتبارسنجی و نتایج، خروجی‌های شبیه‌سازی با روابط تحلیلی یا سناریوهای مرجع مقایسه شده و اثر پارامترهای مهم بر عملکرد سامانه بررسی می‌شود.
	
	\item در صورت ارائهٔ سناریوهای شبکه‌ای، فصل مربوط به شبکه نحوهٔ ترکیب لینک‌های \lr{QKD}، مدیریت استخر کلید، مسیرهای جایگزین و معماری مبتنی بر رلهٔ مورد اعتماد را بررسی می‌کند.
	
	\item در پیوست‌ها، جزئیات نصب، اجرای کد، ساختار فایل‌ها، پیکربندی‌ها و اطلاعات تکمیلی لازم برای بازتولید نتایج آورده می‌شود.
\end{itemize}

در نتیجه، این فصل نقش نقشهٔ راه گزارش را ایفا می‌کند. ابتدا زمینهٔ نظری و انگیزهٔ پژوهش را توضیح می‌دهد، سپس چالش‌های گذار از مدل ایده‌آل به سامانهٔ عملی را روشن می‌سازد، مسئله و دامنهٔ پروژه را مشخص می‌کند، نمایی کلی از فریم‌ورک ارائه می‌دهد، و در نهایت خواننده را برای ورود به فصل‌های تخصصی‌تر آماده می‌کند.
